\documentclass{article}
\usepackage[margin=1.15in]{geometry}
\usepackage{xcolor}
\usepackage{parskip}
\usepackage{fancyhdr}
\usepackage{array}
\usepackage{booktabs}
\usepackage{enumitem}
\usepackage{amsmath}
\usepackage{tabularx}

\pagestyle{fancy}
\fancyhf{}
\cfoot{\thepage}
\rhead{Lecture 7}
\lhead{MA 214: Applied Statistics (Spring 2026)}
\renewcommand{\headrulewidth}{.4mm}

\newcommand{\blankline}[1]{\underline{\hspace{#1}}}

\begin{document}

\noindent\textbf{Name:}\blankline{2.25in}\hfill\textbf{BUID:}\blankline{1.55in}\newline

\begin{center}
 \Large In-Class Worksheet: Model Selection
\end{center}

\subsection*{Scenario}
Today we practice choosing among candidate models. A larger model may fit the observed data better, but model selection asks whether the improvement is worth the added complexity.

\medskip
\noindent\textbf{Goals:} \textbf{(1)} compare $R^2$, adjusted $R^2$, and AIC, \textbf{(2)} explain why complexity penalties are useful, \textbf{(3)} recognize multicollinearity, \textbf{(4)} choose a model using context and diagnostics.

\subsection*{Part 1: Why Not Just Add Everything?}
A researcher is predicting course evaluation score. They are considering the following models.

\begin{center}
\renewcommand{\arraystretch}{1.35}
\begin{tabular}{lcc}
\toprule
\textbf{Model} & \textbf{Adjusted }\boldmath$R^2$ & \textbf{AIC} \\
\midrule
M1: beauty + age & 0.071 & 720.4 \\
M2: M1 + gender & 0.083 & 714.1 \\
M3: M2 + native language & 0.084 & 713.8 \\
M4: M3 + random noise variable & 0.083 & 715.9 \\
\bottomrule
\end{tabular}
\end{center}

\begin{enumerate}[label={\bfseries (\Alph*)},itemsep=.8em]
\item Based only on adjusted $R^2$, which model looks best? \hfill \blankline{1.3in}
\item Based only on AIC, which model looks best? \hfill \blankline{1.3in}
\item Does M4 improve the model? Explain using both adjusted $R^2$ and AIC. \\[3.5em]
\item Which model would you choose for a first report? Briefly explain. \\[4em]
\end{enumerate}

\subsection*{Part 2: $R^2$ vs. Adjusted $R^2$}
Suppose four linear models are fit to the same response and same dataset.

\begin{center}
\renewcommand{\arraystretch}{1.35}
\begin{tabularx}{\textwidth}{lXcc}
\toprule
\textbf{Model} & \textbf{Predictors} & \boldmath$R^2$ & \textbf{Adjusted }\boldmath$R^2$ \\
\midrule
A & $x_1$ & 0.40 & 0.38 \\
B & $x_1+x_2$ & 0.45 & 0.41 \\
C & $x_1+x_2+x_3$ & 0.46 & 0.39 \\
D & $x_1+x_2+x_3+x_4$ & 0.47 & 0.36 \\
\bottomrule
\end{tabularx}
\end{center}

\begin{enumerate}[label={\bfseries (\Alph*)},itemsep=.8em]
\item Why does $R^2$ keep increasing as predictors are added? \\[3em]
\item Which model would adjusted $R^2$ favor? \hfill \blankline{1.2in}
\item In one sentence, explain why adjusted $R^2$ is more useful than $R^2$ for comparing these models. \\[4em]
\end{enumerate}

\subsection*{Part 3: AIC for Logistic Regression}
Two logistic regression models are fit to predict Donner Party survival.

\begin{center}
\renewcommand{\arraystretch}{1.35}
\begin{tabular}{lcc}
\toprule
\textbf{Model} & \textbf{Predictors} & \textbf{AIC} \\
\midrule
Simple logistic & Age & 60.291 \\
Multiple logistic & Age + Sex & 57.256 \\
\bottomrule
\end{tabular}
\end{center}

\begin{enumerate}[label={\bfseries (\Alph*)},itemsep=.8em]
\item Which model is favored by AIC? \hfill \blankline{1.8in}
\item Does lower AIC prove that the model is true? Explain briefly. \\[3em]

\end{enumerate}

\subsection*{Part 4: Multicollinearity}
A model predicts poverty rate using two demographic predictors: \texttt{female\_house} and \texttt{white}. The predictors are strongly associated with each other.

\begin{center}
\renewcommand{\arraystretch}{1.35}
\begin{tabular}{llcc}
\toprule
\textbf{Model} & \textbf{Term} & \textbf{Estimate} & \textbf{p-value} \\
\midrule
One predictor & female\_house & 0.69 & 0.00 \\
Two predictors & female\_house & 0.89 & 0.00 \\
Two predictors & white & 0.04 & 0.29 \\
\bottomrule
\end{tabular}
\end{center}

\begin{enumerate}[label={\bfseries (\Alph*)},itemsep=.8em]
\item What does it mean to interpret the coefficient of \texttt{female\_house} while holding \texttt{white} fixed? \\[3.5em]
\item Why can strong association among predictors make coefficient interpretation difficult? \\[4em]
\item Name one way to notice possible multicollinearity. \hfill \blankline{2.2in}
\end{enumerate}

\subsection*{Wrap-Up}
Complete the sentence: \emph{A good model-selection decision should balance}

\noindent\blankline{6in}. \\[3em]

\end{document}
